\chapter{Das Dualsystem}

In den digitalen Schaltungen mit seinen zwei Bereichen H und L wird das Dualsystem für die Beschreibung mathematischer Zusammenhänge angewendet. Dabei werden nur zwei Zahlen benötigt, nämlich 1 und 0. Eine Verdeutlichung des Dualsystems ist am einfachsten durch eine Gegenüberstellung mit unserem üblichen Dezimalsystem zu erreichen:\index{Dualsystem}\index{Dualsystems}\index{Dezimalsystem}

\begin{table}[htbp]
  \centering
  \begin{tabular}{cc}
    \toprule
    \textbf{Dezimalsystem} & \textbf{Dualsystem} \\
    \midrule
     0 &  0 \\
     1 &  1 \\
     2 &  10 \\
     3 &  11 \\
     4 &  100 \\
     5 &  101 \\
     6 &  110 \\
     7 &  111 \\
     8 &  1000 \\
     9 &  1001 \\
     10 &  1010 \\
     11 &  1011 \\
    \bottomrule
  \end{tabular}
  \par\medskip usw.
\end{table}

Dabei spricht man die duale Zahl 1001 nicht etwa „Eintausendeins“, sondern man spricht die Zahlen einzeln von links beginnend, also „Eins, Null, Null, Eins“.

Durch diese Darstellung können digitale Schaltungen (z.B. Zähler) sehr einfach und durchsichtig dargestellt werden.

